\section{Related work}

\paragraph{From static web tasks to version-aware evaluation.}
Web-agent evaluation is moving from frozen, single-version tasks toward settings that better approximate deployment. Early benchmarks made browser control measurable with reproducible websites, DOM observations, screenshots, and action traces~\citep{zhou2023webarena,deng2023mind2web,koh2024visualwebarena}. Newer benchmarks add live content, deterministic replicas of real websites, long-horizon chores, and user-centered assistance where instructions can be ambiguous or evolve over time~\citep{yuan2024webcanvas,real2025,webchorearena2025,realwebassist2025}. This shift reflects a broader concern that apparent progress can be inflated by contamination, stale benchmarks, and weak agreement with human judgments~\citep{xue2025illusion}. Time-aware evaluation makes the issue especially clear: agents tuned on one web design may fail across historical interface versions~\citep{timewarp2026}. \webcoevo{} follows this version-aware direction, but treats website versions as a design variable rather than only a fixed test distribution.

\paragraph{Website evolution as task-semantic drift.}
Website change is not merely cosmetic. A valid task can keep the same user goal while changing the entry point, readiness timing, labels, workflow checkpoints, route equivalence, or login return behavior. Web-evolution and testing work already separate superficial page differences from semantic structure, dynamic UI state, visual presentation, workflow order, and authentication hazards~\citep{shao2021webevo,mesbah2012crawling,macchi2021imagebased,sanchez2011workflow,jonker2018shepherd}. Recent interface-for-agent work takes the complementary route of exposing semantic component interfaces or structured GUI knowledge so agents depend less on brittle visual shortcuts~\citep{ci4a2026,graphpilot2026}. Our drift labels---access, surface, content, runtime, process, structural, and functional---are therefore an operational language for valid completion-path changes, not a universal taxonomy of web evolution.

\paragraph{Experience loops for self-evolving agents.}
LLM-agent research is also shifting from one-shot prompting toward reusable experience. Reflection and experiential-learning methods convert trajectories into natural-language memories or rules~\citep{shinn2023reflexion,zhao2023expel}. Recent self-evolving-agent work broadens this pattern to memory, tools, policies, architectures, user interaction, and multi-agent feedback~\citep{gao2025selfevolvingagents}. In practical systems, the same trend appears as interaction-derived data, backward construction of instructions, distilled strategic principles, and policy updates from experience rather than static supervision alone~\citep{learnbyinteract2025,wu2025evolver,liang2025sage,comas2025}. \webcoevo{} uses this idea narrowly: the base web agent remains fixed, while paired control/drift traces update a compact cross-version reflection rulebook.

\paragraph{Adaptive environments and agent-environment co-evolution.}
A parallel line of work treats environments as active curricula rather than passive test sets. Unsupervised environment design, replay-guided adversarial design, and compositional environment generation all use the learner's current capability to decide which tasks should be generated or replayed next~\citep{dennis2020paired,jiang2021replayguided,gur2021code}. Newer LLM-agent work extends this idea with synthetic interaction data and generative simulators that produce tasks near the agent's useful difficulty frontier~\citep{learnbyinteract2025,genenv2025}. \webcoevo{} brings this co-evolutionary view to web-agent robustness under stricter constraints: each generated \linkding{} overlay must boot, preserve evaluator semantics, remain manually solvable, and expose a meaningful behavioral failure rather than an infrastructure artifact.

\paragraph{Positioning.}
These lines leave a missing middle layer. Version-aware benchmarks show that web-agent success is unstable across interface change; web-evolution work explains why such change is semantic; self-evolving-agent work shows how experience can become reusable guidance; and adaptive environment generation shows how the next task distribution can target an agent's frontier. \webcoevo{} connects them in one auditable loop: profile failures, generate valid website drift, mine paired trajectories, update compact reflection rules, and test whether the update preserves earlier successes.
